\documentclass[11pt,a4paper]{article}

\usepackage[margin=0.85in]{geometry}
\usepackage[T1]{fontenc}
\usepackage[utf8]{inputenc}
\usepackage{lmodern}
\usepackage{amsmath}
\usepackage{microtype}
\usepackage{booktabs}
\usepackage{array}
\usepackage{graphicx}
\usepackage{enumitem}
\usepackage{xcolor}

\setlist{topsep=0.2em, itemsep=0.2em, parsep=0pt}

\title{\textbf{GabayPoz Recommender System --- Version 1}\\[0.4em]
       \large Demonstration of the Six-Phase Pipeline}
\author{GabayPoz Project --- Team 4 (Model Development)}
\date{\today}

\begin{document}

\maketitle

\begin{abstract}
\noindent
This document is an executable demonstration of the GabayPoz recommender,
version~1. It is generated by running \texttt{demo.py} against the bundled
sample fixtures under \texttt{data/sample/} and emitting a LaTeX fragment
(\texttt{results.tex}) that this wrapper includes verbatim. Each subsection
describes one synthetic student profile, the inputs the recommender
received, the top-three recommendations it produced, and a deep dive on the
rank-1 pick: matched dimensions, applied constraints, scholarship context,
the PIDS Graduate Tracer Study employment outlook, and the hedged
explanation text rendered to the student. Rebuilding the PDF reproduces
this content bit-for-bit; see \texttt{Makefile}.
\end{abstract}

\vspace{0.5em}
\hrule
\vspace{0.5em}

\section*{1. About this demonstration}

GabayPoz is a constraint-aware, explainable college program recommender
for Grade~11/12 students in Pozorrubio, Pangasinan. Two rounds of EDA
established that the binding constraints on tertiary access are not
academic preparation but \emph{geography} (mean barangay-to-university road
distance: 76.9~km / 79.2~min) and \emph{household economics} (a 12.7$\times$
education-spending gradient across FIES income deciles). The recommender
honours those findings by applying hard mobility and economic filters
\emph{before} any program-fit scoring takes place, and by emitting a
structured 20-field explanation alongside every recommendation.

The pipeline runs in six phases:

\begin{enumerate}\setlength\itemsep{0.15em}
  \item \textbf{Student profile.} Q1--Q9 build a six-dimensional
        affinity vector $S \in [0,5]^6$ over (STEM, HEALTH, ARTS,
        BUSINESS, EDUCATION, AGRICULTURE), with weights
        $s_d = 0.50\,\text{internal}_d + 0.40\,\text{aptitude}_d
              + 0.10\,\text{market}_d$. The market component is held
        flat at $1.0$ in version~1 (TDS \S 4.5).
  \item \textbf{Hard filters.} Q11 (commute-time ceiling per
        student-barangay $\to$ university) and Q10 (combined
        tuition-plus-transport burden against an FIES-anchored
        affordability tier) are applied to the program-school
        universe. A three-step relaxation cascade
        (Q11-only, Q10-only, both) precedes any \texttt{NO\_CANDIDATES}
        error.
  \item \textbf{Fit scoring.} A blended
        $\text{base} = 0.70\,\cos(S,P) + 0.30\,(S\cdot P)/150$ favours
        directional alignment in the cold-start regime.
  \item \textbf{Adjustments.} Q12 duration / board-exam soft penalty
        ($\delta = 0.85$ when program tolerance level exceeds the
        student's), then CHED / RDC priority-program boosters
        ($\beta = 1.10\;|\;1.10\;|\;1.18$ capped --- never $1.21$).
  \item \textbf{Diversity-enforced ranking.} Greedy top-3 with a
        $0.15$ score-gap threshold, so when ranks 1 and 2 fall in the
        same dominant dimension, the rank-3 slot opens to a
        different field unless the score gap is large.
  \item \textbf{Explanation.} A 20-field structured JSON with
        matched dimensions, applied constraints, penalties and
        boosters, scholarship context (including the
        \texttt{mobility\_cuts\_scholarship} interaction flag),
        the PIDS GTS employment outlook, the
        \texttt{agriculture\_compound\_flag}, the
        \texttt{low\_confidence\_flag}, and a hedged human-readable
        text. Field-derivation rules and the schema in full are in
        \texttt{docs/tds.md} \S 6.
\end{enumerate}

\noindent
The four profiles below were chosen to exercise each of the
distinctive behaviours: the GTS employment outlook on a HEALTH
recommendation (D1), the mobility-scholarship interaction (D2),
the agriculture compound-disadvantage flag (D3), and the cumulative
constraint-relaxation cascade (D4).

\vspace{0.5em}
\hrule
\vspace{0.5em}

\section*{2. Demonstrative profiles}

\subsection*{D1: Health-oriented STEM strand, ready for board-exam programs}
\noindent Model: \texttt{tds\_recommender\_v1\_2}; Barangay id: \texttt{1}; Q7: \texttt{STEM}; Q10: \texttt{C}; Q11: \texttt{C}; Q12: \texttt{A}.

\paragraph{Top recommendations.}
\begingroup\small\setlength{\tabcolsep}{6pt}
\noindent\begin{tabular}{@{}c p{4.4cm} p{4.4cm} r@{}}
\toprule
Rank & Program & Primary school & Score \\
\midrule
1 & BS Nursing & PSU Lingayen & 0.746 \\
2 & BS Medical Technology & PSU Lingayen & 0.669 \\
3 & Bachelor of Elementary Education & PSU Lingayen & 0.198 \\
\bottomrule
\end{tabular}\endgroup

\paragraph{Rank-1 explanation.}
BS Nursing matches your strongest Health signals. PSU Lingayen is the primary school suggestion because it passed your travel and affordability limits. Local Health field presence is included only as small context, not a job guarantee.


\subsection*{D2: Arts-oriented HUMSS strand, nearby-only travel}
\noindent Model: \texttt{tds\_recommender\_v1\_2}; Barangay id: \texttt{1}; Q7: \texttt{HUMSS}; Q10: \texttt{B}; Q11: \texttt{A}; Q12: \texttt{B}.

\paragraph{Top recommendations.}
\begingroup\small\setlength{\tabcolsep}{6pt}
\noindent\begin{tabular}{@{}c p{4.4cm} p{4.4cm} r@{}}
\toprule
Rank & Program & Primary school & Score \\
\midrule
1 & AB Communication & PSU Lingayen & 0.760 \\
2 & BS Business Administration & PSU Lingayen & 0.328 \\
3 & BS Computer Science & PSU Lingayen & 0.195 \\
\bottomrule
\end{tabular}\endgroup

\paragraph{Rank-1 explanation.}
AB Communication matches your strongest Arts signals. PSU Lingayen is the primary school suggestion because it passed your travel and affordability limits. Local Arts field presence is included only as small context, not a job guarantee.


\subsection*{D3: Agriculture-oriented TVL strand, moderate budget}
\noindent Model: \texttt{tds\_recommender\_v1\_2}; Barangay id: \texttt{3}; Q7: \texttt{TVL}; Q10: \texttt{B}; Q11: \texttt{B}; Q12: \texttt{B}.

\paragraph{Top recommendations.}
\begingroup\small\setlength{\tabcolsep}{6pt}
\noindent\begin{tabular}{@{}c p{4.4cm} p{4.4cm} r@{}}
\toprule
Rank & Program & Primary school & Score \\
\midrule
1 & BS Agriculture & PSU Lingayen & 0.825 \\
2 & BS Civil Engineering & PSU Lingayen & 0.168 \\
3 & BS Medical Technology & PSU Lingayen & 0.159 \\
\bottomrule
\end{tabular}\endgroup

\paragraph{Rank-1 explanation.}
BS Agriculture matches your strongest Agriculture signals. PSU Lingayen is the primary school suggestion because it passed your travel and affordability limits. Local Agriculture field presence is included only as small context, not a job guarantee.


\subsection*{D4: STEM-oriented GAS strand, moderate budget}
\noindent Model: \texttt{tds\_recommender\_v1\_2}; Barangay id: \texttt{2}; Q7: \texttt{GAS}; Q10: \texttt{B}; Q11: \texttt{B}; Q12: \texttt{B}.

\paragraph{Top recommendations.}
\begingroup\small\setlength{\tabcolsep}{6pt}
\noindent\begin{tabular}{@{}c p{4.4cm} p{4.4cm} r@{}}
\toprule
Rank & Program & Primary school & Score \\
\midrule
1 & BS Computer Science & PSU Lingayen & 0.775 \\
2 & BS Civil Engineering & PSU Lingayen & 0.672 \\
3 & BS Medical Technology & PSU Lingayen & 0.393 \\
\bottomrule
\end{tabular}\endgroup

\paragraph{Rank-1 explanation.}
BS Computer Science matches your strongest STEM signals. PSU Lingayen is the primary school suggestion because it passed your travel and affordability limits. Local STEM field presence is included only as small context, not a job guarantee.


\vspace{1em}
\hrule
\vspace{0.5em}

\section*{3. How to reproduce this PDF}

\begin{verbatim}
cd demo
make            # runs demo.py then pdflatex twice
\end{verbatim}

\noindent
\texttt{make clean} removes the LaTeX intermediates and the generated
\texttt{results.tex}; \texttt{make distclean} additionally removes
\texttt{demo.pdf}. Re-running from scratch reproduces the PDF
deterministically: the recommender is fully stateless, the seven
sample CSVs are version-controlled, and tie-breaking is documented
in TDS \S 5 Phase~5.

\vspace{0.5em}
\hrule
\vspace{0.5em}

\section*{4. References}

\begin{itemize}\setlength\itemsep{0.15em}
  \item \textbf{TDS:} \texttt{docs/tds.md} --- the authoritative
        specification this implementation traces to.
  \item \textbf{EDA v1:} \texttt{docs/eda\_v1.pdf} --- Pangasinan
        education, employment, and access profile.
  \item \textbf{EDA v2:} \texttt{docs/eda\_v2.pdf} --- Pozorrubio
        program access, scholarship analysis, and employment signals.
  \item Burke, R.\ (2002). Hybrid recommender systems.
        \emph{User Modeling and User-Adapted Interaction}, 12(4).
  \item Liu, T.-Y.\ (2009). Learning to rank for information
        retrieval. \emph{Foundations and Trends in IR}, 3(3).
  \item Philippine Institute for Development Studies (2019).
        \emph{Tracer study of CHED-scholarship and non-scholarship
        graduates} (Discussion Paper No.\ 2019-26).
\end{itemize}

\end{document}
